\documentclass[11pt]{article}

\usepackage[margin=1in]{geometry}
\usepackage{amsmath,amssymb,amsthm,mathtools}
\usepackage{graphicx}
\usepackage{microtype}
\usepackage{xcolor}
\usepackage[hidelinks]{hyperref}
\usepackage{enumitem}

\definecolor{deepblue}{RGB}{28,62,103}
\newtheorem{theorem}{Theorem}
\newtheorem{lemma}{Lemma}
\newtheorem{remark}{Remark}
\newcommand{\C}{\mathbb C}
\newcommand{\N}{\mathbb N}
\newcommand{\e}{\mathrm e}

\title{A Block-Clustered Generator with a Bounded Non-Riesz\\
Basis of Eigenvectors}
\author{Candidate full solution packet for arXiv:1405.2731}
\date{9 August 2026}

\begin{document}
\maketitle

\begin{center}
\fbox{\begin{minipage}{0.9\textwidth}
\textbf{Status.} Candidate full solution, likely valid, pending human review.
The construction gives an affirmative answer and also produces simple purely
imaginary eigenvalues and a linearly bounded group.
\end{minipage}}
\end{center}

\section{Source question}

Sklyar and Marchenko ask in the concluding remarks of
\cite[p.~32]{SklyarMarchenko}:
\begin{quote}
Is it possible to construct the unbounded generator of a $C_0$-group with
bounded non-Riesz basis family of eigenvectors?
\end{quote}
Here ``bounded non-Riesz basis'' means a Schauder basis whose vector norms are
bounded above and away from zero, but which is not a Riesz basis; see also the
authors' later formulation in arXiv:1809.03079.  The source passage is shown
below.

\begin{figure}[ht]
  \centering
  \includegraphics[width=0.98\textwidth]{figures/open_problem_crop.png}
  \caption{The complete question and the semigroup comparison on page 32 of
  the source paper.}
\end{figure}

\section{Idea of the proof}

Start from one normalized conditional Schauder basis in a Hilbert space and
look at longer and longer finite prefixes.  Their basis constants remain
uniform, but their Hilbertian condition numbers become arbitrarily large.
Place copies of these prefixes in mutually orthogonal finite-dimensional
blocks.  The concatenation is still a normalized Schauder basis, while its
block condition numbers force it not to be Riesz.

An arbitrary diagonal phase multiplier can be large on an ill-conditioned
block.  The key is to give the vectors in that block almost the same
frequency.  The common phase is an isometry and disappears from the norm
estimate; only the within-block frequency diameter is amplified by the bad
condition number.  Choosing that diameter inversely proportional to the
condition number yields a uniform group bound.  Meanwhile the common
frequencies of successive blocks tend to infinity, so the generator is
unbounded.  This simultaneously reconciles strong continuity, unbounded
spectrum, and conditional basis geometry.

\section{The result}

\begin{theorem}\label{thm:main}
There exist a separable complex Hilbert space $H$, a strongly continuous group
$(T(t))_{t\in\mathbb R}$ on $H$, and its unbounded generator $A$ such that:
\begin{enumerate}[label=\textup{(\roman*)}]
\item $A$ has pairwise distinct purely imaginary eigenvalues;
\item corresponding normalized eigenvectors form a Schauder basis of $H$ but
      not a Riesz basis; and
\item
\[
             \|T(t)\|\le 1+\frac{|t|}{2}\qquad(t\in\mathbb R).
\]
\end{enumerate}
In particular, the answer to the source question is affirmative.
\end{theorem}

\section{Proof}

The existence of a normalized conditional Schauder basis in a separable
Hilbert space is classical.  For example, the source paper
\cite[p.~7]{SklyarMarchenko} recalls Babenko's explicit bounded non-Riesz basis
in $L_2(-\pi,\pi)$ \cite{Babenko}.  Fix such a basis $(u_j)_{j\ge1}$ in a
Hilbert space $H_0$, and let
\[
 K=\sup_{r\ge1}\left\|\sum_{j=1}^r u_j^*\otimes u_j\right\|<\infty
\]
be its Schauder basis constant.

For $N\ge1$, put $E_N=\operatorname{span}\{u_1,\ldots,u_N\}$ and let
\[
 S_N:\ell_2^N\longrightarrow E_N,\qquad S_N e_j=u_j,
 \qquad \kappa_N=\|S_N\|\,\|S_N^{-1}\|.
\]

\begin{lemma}\label{lem:condition}
The sequence $(\kappa_N)$ is unbounded.
\end{lemma}

\begin{proof}
Normalization gives $\|S_N\|\ge1$ and $\|S_N^{-1}\|\ge1$.  If
$\sup_N\kappa_N<\infty$, both factors would be uniformly bounded.  Hence a
constant $C$ would satisfy
\[
 C^{-1}\Big(\sum_{j=1}^N|c_j|^2\Big)^{1/2}
 \le \left\|\sum_{j=1}^N c_j u_j\right\|
 \le C\Big(\sum_{j=1}^N|c_j|^2\Big)^{1/2}
\]
for every $N$ and every scalar family $(c_j)_{j=1}^N$.  These are precisely
uniform Riesz inequalities, contradicting that $(u_j)$ is non-Riesz.
\end{proof}

Choose strictly increasing $N_m$ such that
$\kappa_m:=\kappa_{N_m}\to\infty$.  Let $H_m$ be an orthogonal copy of
$E_{N_m}$, write $(v_{m,j})_{j=1}^{N_m}$ for the copied basis, and retain the
notation
\[
 S_m:\ell_2^{N_m}\longrightarrow H_m,\qquad S_m e_j=v_{m,j}.
\]
Set
\[
                  H=\left(\bigoplus_{m=1}^{\infty}H_m\right)_{\ell_2}
\]
and order the vectors $(v_{m,j})$ block by block.

\begin{lemma}\label{lem:basis}
The blockwise concatenation $(v_{m,j})$ is a normalized Schauder basis of $H$
with basis constant at most $\max(1,K)$, and it is not a Riesz basis.
\end{lemma}

\begin{proof}
Let $P_{m,r}$ be the projection in $H_m$ onto the first $r$ copied basis
vectors.  It is the restriction of the corresponding partial-sum projection
of $(u_j)$, so $\|P_{m,r}\|\le K$.  A global partial-sum projection retains
all blocks before some $m$, applies $P_{m,r}$ in block $m$, and annihilates
later blocks.  Orthogonality gives, for $x=(x_q)_q$,
\[
 \left\|\big(x_1,\ldots,x_{m-1},P_{m,r}x_m,0,\ldots\big)\right\|^2
 \le \max(1,K^2)\sum_q\|x_q\|^2.
\]
Uniform boundedness of these projections and density of the finite block
vectors prove the Schauder basis assertion.

If the concatenation had Riesz bounds independent of the coefficients, their
restriction to block $m$ would uniformly bound both $\|S_m\|$ and
$\|S_m^{-1}\|$.  This would bound $\kappa_m$, a contradiction.
\end{proof}

We now assign frequencies.  For $m\ge1$ put
\[
 a_m=4^m,
 \qquad
 \delta_m=\frac{2^{-m}}{\kappa_m N_m},
 \qquad
 \beta_{m,j}=a_m+\delta_m j\quad(1\le j\le N_m).
\]
Since $0<\delta_mN_m\le2^{-m}<1$, the frequencies in block $m$ lie in
$(4^m,4^m+1)$.  Thus all $\beta_{m,j}$ are positive and pairwise distinct.
On $H_m$ define
\[
 T_m(t)=S_m\,\operatorname{diag}
       (\e^{i\beta_{m,1}t},\ldots,\e^{i\beta_{m,N_m}t})\,S_m^{-1}.
\]
After removing the common scalar phase $\e^{ia_mt}$,
\begin{align*}
 \|T_m(t)\|
 &\le 1+\|S_m\|\,\|S_m^{-1}\|
        \max_{1\le j\le N_m}|\e^{i\delta_mjt}-1|\\
 &\le 1+\kappa_m\delta_mN_m|t|
  =1+2^{-m}|t|.                                      \tag{1}\label{eq:block}
\end{align*}
Here we used $|\e^{is}-1|\le |s|$ for real $s$.

Define the orthogonal direct sum
\[
                       T(t)=\bigoplus_{m=1}^{\infty}T_m(t).
\]
Estimate \eqref{eq:block} shows that this is a bounded operator and
\[
                    \|T(t)\|=\sup_m\|T_m(t)\|
                    \le 1+\frac{|t|}{2}.              \tag{2}\label{eq:global}
\]
The group law holds in every block and hence globally.  On the dense subspace
of vectors supported in finitely many blocks, $T(t)x\to x$ as $t\to0$ by
finite-dimensional continuity.  Estimate \eqref{eq:global} is uniform for
$|t|\le1$, so a density argument extends this convergence to every $x\in H$.
Thus $(T(t))_{t\in\mathbb R}$ is a $C_0$-group.

For completeness, we identify its generator without invoking a direct-sum
generator theorem.  Put
\[
 A_m=S_m\,\operatorname{diag}
       (i\beta_{m,1},\ldots,i\beta_{m,N_m})\,S_m^{-1}
\]
and define
\[
 \begin{split}
 D(A)&=\{x=(x_m)_m\in H:(A_mx_m)_m\in H\},\\
 Ax&=(A_mx_m)_m.
 \end{split}                                           \tag{3}\label{eq:domain}
\]
If $x\in D(A)$ and $y=Ax$, the finite-dimensional identities
\[
 T_m(t)x_m-x_m=\int_0^tT_m(s)A_mx_m\,ds
\]
assemble, using strong continuity and the local bound, into
\[
                  T(t)x-x=\int_0^tT(s)y\,ds.
\]
Dividing by $t$ and letting $t\to0$ proves that $x$ belongs to the generator
domain and that the generator sends $x$ to $Ax$.  Conversely, if $x$ belongs
to the generator domain, applying the bounded orthogonal projection onto
$H_m$ to its difference quotient gives $(Ax)_m=A_mx_m$.  The generator value
lies in $H$, so $(A_mx_m)_m\in H$.  This proves \eqref{eq:domain} is exactly
the generator domain.

Finally,
\[
             Av_{m,j}=i\beta_{m,j}v_{m,j}.
\]
The eigenvalues are distinct and purely imaginary, and Lemma
\ref{lem:basis} gives the required normalized non-Riesz Schauder basis of
eigenvectors.  Moreover,
\[
                 \|Av_{m,1}\|=\beta_{m,1}\ge4^m
\]
while $\|v_{m,1}\|=1$, so $A$ is unbounded.  This proves the theorem.
\qed

\section{Verification report}

The proof is analytic and does not depend on finite computation.  The checks
most likely to expose a defect were carried out separately:
\begin{enumerate}[leftmargin=*,label=\arabic*.]
\item Uniform partial-sum bounds survive concatenation because the blocks are
orthogonal; Lemma \ref{lem:basis} gives the exact estimate.
\item The basis cannot accidentally become Riesz, because Riesz bounds would
hold on every block and contradict $\kappa_m\to\infty$.
\item The only amplification by the ill-conditioned basis is
$\kappa_m\delta_mN_m$, which equals $2^{-m}$ exactly.  This proves boundedness
for every $t$ and supplies the uniform local bound needed for strong
continuity.
\item Formula \eqref{eq:domain} was verified in both directions through the
integral identity and block projections.  Unboundedness is witnessed on unit
eigenvectors, not inferred merely from the written eigenvalue formula.
\end{enumerate}
No unproved lemma or computational dependency remains.  A longer checklist is
included as \texttt{verification\_report.md} in the packet.

\section{Novelty check, limitations, and review recommendation}

A bounded search was performed on 9 August 2026 in the run's local arXiv
corpus and through web/arXiv search.  Queries included the exact sentence
``unbounded generator ... bounded non-Riesz basis'', the authors' names, and
combinations of ``conditional Schauder basis'', ``eigenvectors'', and
``$C_0$-group''.  The search also inspected the authors' later arXiv:1809.03079
and the related 2021 work on resolvents and polynomial bounds.  Those works
treat complete minimal eigenvector families that do not form a Schauder basis;
no later source explicitly answering the bounded Schauder non-Riesz question
was found.  This is a bounded novelty check, not a certification of priority.

The theorem answers the construction question but does not address Conjecture
20 of the source paper, which imposes a finite-union uniform-gap condition on
the spectrum.  Indeed, the close frequency clusters here deliberately violate
that condition.  Human review should focus on Lemma \ref{lem:basis} and
estimate \eqref{eq:block}; the rest is a standard but fully written direct-sum
group argument.

\begin{thebibliography}{9}

\bibitem{SklyarMarchenko}
G.~M. Sklyar and V.~Marchenko,
\emph{Hardy inequality and the construction of infinitesimal operators with
non-basis family of eigenvectors},
J. Funct. Anal. \textbf{272} (2017), 1017--1043;
arXiv:1405.2731v4.

\bibitem{Babenko}
K.~I. Babenko,
\emph{On conjugate functions},
Dokl. Akad. Nauk \textbf{62} (1948), no.~2, 157--160 (Russian),
MR0027093.

\end{thebibliography}

\end{document}
